\ \\
\Def \textit{Комплексное число z} -- это выражение вида $z = a+ib$, где $a$ и $b$ -- числа из $\R$, а $i$ -- мнимая единица. По определению
$i^2 = -1$. Число $a$ называют вещественной частью $z$ (пишется $a = Re (z))$, а число $b$ -- мнимой частью $z$ (пишется
$b = Im(z))$. Комплексные числа можно складывать и умножать, «раскрывая скобки и приводя подобные». Множество
комплексных чисел обозначают буквой $\mathbb{C}$.
\\
\\
\Def \textit{Сложение} комплексных чисел определяется как $a_1 + i b_1 + a_2 + i b_2 = (a_1 + a_2) + i(b_1 + b_2)$
\\
\Def \textit{Умножение} комплексных чисел определяется как $(a_1 + i b_1)*( a_2 + i b_2) = a_1a_2 - b_1b_2 + i(a_1b_2 + a_2b_1)$
\\
\Def
Каждому комплексному числу $z = a + bi$ сопоставим точку $(a, b)$ и вектор $(a, b)$. Длина этого вектора называется
\textit{модулем числа z} и обозначается $|z|$
\\
\Def Пусть $z \neq 0$.
Угол (в радианах), отсчитанный против часовой стрелки от вектора (1, 0) до вектора (a, b), называется \textit{аргументом числа z} и обозначается $Arg(z)$. Аргумент определен с точностью до прибавления
числа вида $2\pi n$, где $n \in Z$.
\\
\Def \textit{Тригонометрическая форма записи}. Для любого ненулевого комплексного числа $z$ имеет место равенство
$z = r(cos \phi + i sin \phi)$, где $r = |z|$, $\phi = Arg (z)$. \\
Для комплексного числа $z = r(cos \phi + i sin \phi)$ и натурального числа $n \in \N$ выполнена \textit{формула Муавра}: \begin{center}
    $z^n = r^n(cos( n  \phi) + i sin ( n \phi))$.
\end{center}
\Proof $z^n = (re^{i\phi})^n = r^n e^{i\phi n} = r^n(cos( n 
\phi) + i sin ( n \phi))$
\EndProof \\
\Def Для комплексного числа $z = a+ib$, где $a, b \in \mathbb{R}$ число $z = a - bi$ называется \textit{комплексно-сопряжённым к z}
\\
\\
\textbf{Утверждение:} \\
Выполнены следующие равенства:
\begin{itemize}
    \item[1] $|z|^2 = z\overline{z}$
    \item[2] $ \overline{z + w} = \overline{z} + \overline{w}$
     \item[3] $\overline{z} * \overline{w} = \overline{z w}$
\end{itemize}

\Def \textit{Корнем из 1 степени n} называется такое комплексное число z, что $z^n = 1$.